\documentclass[11pt]{article}
\usepackage[margin=1in]{geometry}
\usepackage{amsmath,amssymb,amsthm,hyperref}
\hypersetup{colorlinks=true,linkcolor=blue,citecolor=blue,urlcolor=blue,
 pdftitle={Valuation balance and local lifting for cyclic n-roots},pdfauthor={Jeff Kline}}
\newtheorem{theorem}{Theorem}[section]
\newtheorem{proposition}[theorem]{Proposition}
\newtheorem{corollary}[theorem]{Corollary}
\theoremstyle{remark}\newtheorem{remark}[theorem]{Remark}
\DeclareMathOperator{\supp}{supp}
\DeclareMathOperator{\diag}{diag}
\DeclareMathOperator{\rank}{rank}
\newcommand{\C}{\mathbb C}\newcommand{\Q}{\mathbb Q}
\newcommand{\F}{\mathbb F}\newcommand{\Z}{\mathbb Z}
\title{Valuation balance and local lifting for cyclic \texorpdfstring{$n$}{n}-roots}
\author{Jeff Kline}
\date{Draft of version 0.1.0 --- October 26, 2026 --- not released}
\begin{document}
\maketitle
\begin{abstract}
In its usual complex form, the Bj\"orck--Saffari conjecture asks whether
the cyclic $n$-roots are finite in number whenever $n$ is square-free. This
is proved for prime $n$; among composite square-free orders, the sources we read
settle only $n=6$ and $n=10$, by computer algebra. Haagerup's prime-order proof starts from the fact that an infinite
solution set yields disjoint supports in time and frequency, and Zhou showed
that such supports exist at every composite order. We show that the Fourier
equations force more. Along any Laurent-series solution, the multiset of
coordinate valuations equals the multiset of valuations of its Fourier
transform. This is a consequence of a residue principle for master
functions, and we claim no priority for it. It gives exact complementary
supports of equal size at every proper valuation gap. For a one-dimensional
endpoint kernel and a nondegenerate projective normal critical point, we
reduce local lifting to one analytic function and give its first two scalar
conditions, $\kappa=0$ and $\mu=0$. At order $26$, one fixed support pair
admits nondegenerate normal critical points with $\kappa\ne0$ and others with
$\kappa=0$; an affine reflection forces the latter cancellation exactly. At
two such points a rigorous $\pi$-adic computation shows $\mu\ne0$, so no
solution branch leaves them in the prescribed directions. Thus $\kappa$ alone
cannot exclude every endpoint, while $\mu$ excludes these. We prove no bound
on the number of conditions needed in general, and no finiteness theorem for
any composite order.
\end{abstract}

\section{The question and the contribution}
A cyclic $n$-root is a vector $a\in\C^n$ satisfying
\[
 \sum_{j=0}^{n-1}a_j a_{j+1}\cdots a_{j+r-1}=0\quad(1\le r<n),
 \qquad \prod_{j=0}^{n-1}a_j=1,
\]
with indices modulo $n$. If $m^2$ divides $n$ for some $m>1$, there are
infinitely many cyclic $n$-roots (Backelin \cite{Backelin}; Bj\"orck and
Saffari for the unimodular solutions \cite{BS95}). Bj\"orck and Saffari
conjectured that for square-free $n$ the unimodular solutions with leading
entry $1$ are finite in number \cite[Conj.~1.1]{FR}; the form for all complex
solutions is the standard restatement \cite{Zhou}. Haagerup proved
finiteness, and a count with multiplicity, for prime order \cite{Haagerup}.
For composite square-free orders, the sources we read establish finiteness
only for $n=6$ and $n=10$, by Gr\"obner-basis computation; at $n=14$ a count
of isolated roots is reported, but finiteness is not established there
(see the table and its footnote in \cite{FR}). The question for $2p$, with $p$ an arbitrary
odd prime, motivates this note but is not resolved here.

Put $\zeta=e^{2\pi i/n}$ and $F=(\zeta^{ij})_{0\le i,j<n}$.
We use the unnormalized Fourier matrix, bilinear transpose, and
coordinatewise reciprocals. Work where all coordinates of $x$ and
$f=Fx$ are nonzero. The equations are
\begin{equation}\label{eq:critical}
 E(x):=F(Fx)^{-1}-x^{-1}=0.
\end{equation}
After fixing $x_0=1$, these are the Fourier equations for cyclic roots:
set $a_j=x_{j+1}/x_j$. Indeed, $F^{-1}(x^{-1})=(Fx)^{-1}$ says that
the cyclic correlations $\sum_j x_{j+r}/x_j$ vanish for $r\ne0$.
This correspondence is credited to the existing literature,
particularly \cite{Haagerup}.

Haagerup's limiting-support argument produces disjoint supports in time
and frequency from an infinite solution set. Zhou proves that such pairs
exist for every composite order \cite{Zhou}. Our valuation statement
retains multiplicities and all the gaps, yielding exact complements of
equal size. These conditions are necessary, not sufficient for a branch.

We then express the normal equations at a boundary point as the critical
equations for a product of linear forms. Such master functions are standard;
see \cite{CDFV}. An implicit-function reduction supplies the concrete scalars
$\kappa,\mu$. Puiseux expansions of cyclic $n$-roots were computed earlier
from initial forms by Adrovic and Verschelde, who test whether a second term
exists, starting from a regular isolated root of the initial system \cite[Sec.~5.1]{AV};
$\kappa$ and $\mu$ are compatibility conditions of this kind, written
explicitly for nondegenerate normal critical points at a valuation gap. The
order-$26$ example shows both nonvanishing and symmetry-forced vanishing of
the first scalar; at two of the symmetric points a rigorous $\pi$-adic
computation shows $\mu\ne0$, so no branch leaves them in the prescribed
directions. The reflection used there is, up to a scalar factor, a
composite of the shifts, modulations and dilations that F\"uhr and Rzeszotnik
list for unimodular solutions \cite[Sec.~10.3]{FR}; the same formulas act on
all solutions. The Fourier symmetries, master-function framework, and
implicit-function method are not claimed as new.

Theorem~\ref{thm:valuations} is a consequence of a general residue
principle: when a formal curve of critical points of a master function
$\prod f_i^{\alpha_i}$ tends to boundary divisors $E$ of a smooth
simple normal-crossings compactification, the weights satisfy
$\sum_i\operatorname{ord}_E(f_i)\alpha_i=0$ for each of them
\cite[Prop.~2.10]{SvdV}. Apply it to the $2n$ forms $x_i$ and $(Fx)_k$, with
weights $+1$ on $(Fx)_k$ and $-1$ on $x_i$ (total weight $0$, so the master
function lives on projective space), on the maximal projective wonderful
model of their arrangement \cite[\S\S3--4]{DP}; see also
\cite[Thm.~2.5, \S\S4.1--4.2]{Fei} and, for its fan,
\cite[Prop.~3.6, Thm.~4.1]{FS}. Its boundary divisors are indexed by the
nonempty proper flats. Each set $\{\ell:\nu(\ell)>c\}$ is a flat, since a
form in the span of forms of order $>c$ has order $>c$, and the solution
curve meets its divisor, along which $\ell$ vanishes to order $1$ exactly when
$\ell$ lies in the flat. The residue condition there says that the flat
contains as many forms $(Fx)_k$ as forms $x_i$, which is the equality of the
counts at threshold $c$. We give a direct proof below and claim no priority
for Theorem~\ref{thm:valuations}. For cyclic $n$-roots, this note adds
the exact supports at every gap, the scalars $\kappa,\mu$, and the order-$26$
example; a bounded search, recorded with the note, found
none of them stated for cyclic $n$-roots.

\section{Valuations and every gap}
Write $\nu(h)$ for the order of a nonzero Laurent series at $t=0$.
The next result uses symmetry and invertibility rather than Fourier entries.

\begin{theorem}[Valuation multiset]\label{thm:valuations}
Let $F$ be an invertible symmetric constant matrix over $\C$. Suppose
$x\in\C((t))^n$ satisfies \eqref{eq:critical} and all coordinates of
$x$ and $Fx$ are nonzero. Then
\[
 \{\nu(x_i):0\le i<n\}=\{\nu((Fx)_k):0\le k<n\}
\]
as multisets. This also holds for Puiseux series after a finite change
of parameter.
\end{theorem}
\begin{proof}
Let $D=t\,d/dt$, $y=x^{-1}$, and $f=Fx$. Since $F^{-1}y=f^{-1}$,
for every integer $r\ge1$ we have
\[
 \sum_k\frac{D^rf_k}{f_k}
 =(FD^rx)^TF^{-1}y=(D^rx)^Ty
 =\sum_i\frac{D^rx_i}{x_i}.
\]
If $h=ct^a(1+O(t))$, then $D^rh/h=a^r+O(t)$. Constant terms give
equal power sums of the valuation multisets. Newton's identities for the
first $n$ sums determine a multiset of $n$ numbers. A common denominator
of Puiseux exponents reduces that case to Laurent series.
\end{proof}

\begin{corollary}[Exact complementarity at every gap]\label{cor:gaps}
Under Theorem~\ref{thm:valuations}, choose a real $c$ which is not a
coordinate valuation, and put
\[
 I=\{i:\nu(x_i)<c\},\qquad J=\{k:\nu((Fx)_k)<c\}.
\]
If the cut is proper, then $0<|I|=|J|<n$, and constant vectors $u,v$ exist with
\[
 \supp u=I,\quad\supp(Fu)=J,\qquad
 \supp v=I^c,\quad\supp(F^{-1}v)=J^c.
\]
\end{corollary}
\begin{proof}
The cardinalities follow from the multiset theorem. Each coefficient vector
$x^{[r]}$ for $r<c$ is supported in $I$ with transform supported in $J$.
There are finitely many such nonzero vectors. Every coordinate of $I$,
and every transform coordinate of $J$, occurs in at least one.
A generic constant linear combination has exactly those supports,
since bad combinations form a finite union of proper hyperplanes.
Apply the same argument to $y=x^{-1}$ and $F^{-1}y=f^{-1}$ at orders
less than $-c$ to obtain the complements.
\end{proof}

\begin{remark}
An infinite normalized cyclic-root variety contains an algebraic curve.
On its smooth projective model a nonconstant coordinate has a pole.
The Laurent expansion there, with $x_0=1$, has a nonconstant valuation
vector and a proper gap. Thus the corollary gives necessary data for an
infinite solution set, without asserting that arbitrary such data occur
along a solution curve.
\end{remark}

\section{The normal equations and two scalar conditions}
Continue with any invertible symmetric $F$. Fix $|I|=|J|=m$, put $q=n-m$,
and suppose
\begin{equation}\label{eq:endpoint}
 \supp u=I,\quad\supp f=J,\quad f=Fu,\quad
 A=F[J^c,I],\quad\rank A=m-1.
\end{equation}
The right kernel is the line $u_I$. Let the columns of $G\in\C^{q\times d}$
span $\ker A^T$, where $d=q-m+1$. Set
\[
 V=F[I^c,J^c]G,\quad g_{J^c}=Gz,\quad g_J=0,\qquad v=Fg.
\]
Work where all entries of $Gz$ and $Vz=v_{I^c}$ are nonzero. Define
\begin{equation}\label{eq:psi}
 \Psi(z)=\sum_{k\in J^c}\log(Gz)_k-\sum_{i\in I^c}\log(Vz)_i.
\end{equation}
Only logarithmic derivatives are used. The equal factor counts make
this a projective critical-point problem.

\begin{proposition}\label{prop:normal}
A vector $w$ with
\begin{equation}\label{eq:w}
 w_{I^c}=v_{I^c}^{-1},\qquad (Fw)_{J^c}=g_{J^c}^{-1}
\end{equation}
exists exactly when $\nabla\Psi(z)=0$, and is unique modulo $u$.
Set
\[
 H=V^T\diag(v_{I^c}^{-2})V-G^T\diag(g_{J^c}^{-2})G,\qquad
 L=\diag(v^2)-F\diag(g^2)F.
\]
At a critical point, $\dim\ker L=1+\dim\ker H$.
If $\rank H=d-1$, then $\ker L=\operatorname{span}\{u,w\}$.
\end{proposition}
\begin{proof}
Solving for $w_I$ is an $A$-system with complete left-kernel conditions
$G^Tg^{-1}-V^Tv^{-1}=0$, precisely $\nabla\Psi=0$.
Let $Nr=(r_{I^c},(Fr)_{J^c})$ and
$D=\diag(v_{I^c}^2,-g_{J^c}^2)$. Then $L=N^TDN$, $\ker N=\C u$.
The columns of $\Lambda=(V^T,-G^T)^T$ span $\ker N^T$; their number is
$2q-(n-1)=d$. Also $\Lambda^TD^{-1}\Lambda=H$.
Indeed $\dim\ker N^T=2q-\rank N=2q-(n-1)=d$, and the $d$ columns of
$\Lambda$ are independent because those of $G$ are.
Now $Lr=0$ means $DNr=\Lambda b$. Such $r$ exists exactly when $Hb=0$,
and is unique modulo $u$. Finally $Lu=0$, and \eqref{eq:w} gives
$DNw=(v_{I^c},-g_{J^c})=\Lambda z$, so $Lw=N^T\Lambda z=0$; $u,w$ are
independent since $w$ is nonzero on $I^c$.
\end{proof}

Fix a critical point $z$ of $\Psi$ and assume projective nondegeneracy:
$\rank H=d-1$. Put
\begin{equation}\label{eq:kappa}
 h=F(1/f\text{ on }J)-(1/u\text{ on }I),\qquad\kappa=w^Th,
\end{equation}
with zero entries elsewhere. Since $u^Th=|J|-|I|=0$, $\kappa$ is
independent of the choice of $w$. Symmetry of $L$ gives
\begin{equation}\label{eq:solvable}
 Ls=-h\text{ is solvable}\quad\Longleftrightarrow\quad\kappa=0.
\end{equation}
When this holds, define
\begin{equation}\label{eq:mu}
 \mu=\sum_{j\in J}\left(\frac{(Fw)_j}{f_j}\right)^2
       -\sum_{i\in I}\left(\frac{w_i}{u_i}\right)^2+s^TLs.
\end{equation}
Changing $s$ by $\ker L$ has no effect. Replacing $w$ by $w+cu$ changes
the two sums by $2c\kappa+c^2(|J|-|I|)=0$.
Under $u\mapsto au$, $g\mapsto bg$, we have
$\kappa\mapsto(ab)^{-1}\kappa$ and, when defined,
$\mu\mapsto(ab)^{-2}\mu$. Their vanishing is independent of scale.

\begin{theorem}[Local scalar reduction]\label{thm:local}
At such a critical point, under \eqref{eq:endpoint}, the nonzero-coordinate
conditions, and projective nondegeneracy, there is a convergent effective function
$Q_{\rm eff}(\tau)$ with $Q_{\rm eff}'(0)=\kappa$. If $\kappa=0$, then
\[
 Q_{\rm eff}'(\tau)=-\mu\tau+O(\tau^2).
\]
A formal branch with boundary direction $u$ and first transverse direction
$w$ exists exactly when $Q_{\rm eff}'\equiv0$. Here such a branch is a
solution $x(t)$ of \eqref{eq:critical} over $\C((t))$, with all
coordinates of $x$ and $Fx$ nonzero, which after multiplication by a
scalar series satisfies $\ell_u(x)=1$ and $x=u+\tau(t)w+\xi(t)$, where
$\ell_u,\ell_w,T$ are as in the proof, $\xi\in T[[t]]$, $\tau\in t\C[[t]]$
is nonzero, and $\nu(\xi)>\nu(\tau)$.
\end{theorem}
\begin{proof}
Choose linear forms $\ell_u,\ell_w$ dual to $u,w$ on their span and put
$T=\ker\ell_u\cap\ker\ell_w$. In the slice $\ell_u(x)=1$, write
$x=u+\tau(w+\eta)$, $\eta\in T$. The rational function
$\mathcal R(x)=\prod_k(Fx)_k/\prod_i x_i$ has logarithmic gradient $E(x)$.
Cancelling its $\tau^q$ factors gives, up to a constant, the logarithm
\begin{align*}
 Q(\tau,\eta)&=Q_N(\eta)+Q_R(\tau,\eta),\\
 Q_N&=\sum_{j\in J^c}\log\left(1+\frac{(F\eta)_j}{(Fw)_j}\right)
       -\sum_{i\in I^c}\log\left(1+\frac{\eta_i}{w_i}\right),\\
 Q_R&=\sum_{j\in J}\log\left(1+\tau\frac{(Fw+F\eta)_j}{f_j}\right)
       -\sum_{i\in I}\log\left(1+\tau\frac{w_i+\eta_i}{u_i}\right).
\end{align*}
At $(0,0)$ the $\eta$-gradient is zero, by the normal system \eqref{eq:w},
and its Hessian is $L|_T$.
A vector in its radical pairs to zero with $T$ and with $u,w$, so lies
in $\ker L\cap T=0$. The formal and analytic implicit-function theorems
give the same unique $\eta(\tau)\in\tau T[[\tau]]$ with
$\partial_\eta Q=0$. Define $Q_{\rm eff}(\tau)=Q(\tau,\eta(\tau))$.

For $\tau\ne0$ these coordinates parametrize the slice near the specified
direction. Homogeneity of $\mathcal R$ makes criticality on the slice
equivalent to $E(x)=0$. Given a branch as in the statement, put
$\eta_b=\xi/\tau\in tT[[t]]$, so that $x=u+\tau(w+\eta_b)$ and both
partial derivatives of $Q$ vanish at $(\tau(t),\eta_b(t))$. Write
$\partial_\eta Q(\tau,\eta)=L|_T\eta+\beta(\tau)+O(\tau\eta)+O(\eta^2)$ with
$\beta(0)=0$. Since $\tau(t)\in t\C[[t]]$ and $L|_T$ is invertible, the
coefficient of $t^k$ in $\partial_\eta Q(\tau(t),\eta(t))=0$ determines the
$t^k$ coefficient of $\eta$ from the lower ones; so this equation has a
unique solution in $tT[[t]]$, and $\eta_b=\eta\circ\tau$. Then
$Q_{\rm eff}'(\tau(t))=\partial_\tau Q+\partial_\eta Q\cdot\eta'=0$, and a
nonconstant $\tau(t)$ cannot annihilate a nonzero power series in $\tau$.
A branch therefore implies
$Q_{\rm eff}'\equiv0$. Conversely that identity supplies the convergent
branch $u+\tau(w+\eta(\tau))$, with all required coordinates nonzero
for small nonzero $\tau$.

Differentiation gives $Q_\tau(0,0)=w^Th=\kappa$.
If $\kappa=0$, let $s=\eta'(0)$. The transverse equation gives
$(Ls+h)|_T=0$; it also vanishes on $u,w$, so $Ls=-h$. Finally,
\[
 Q_{\rm eff}''(0)=Q_{\tau\tau}+2h^Ts+s^TLs
 =-\left\{\sum_J((Fw)_j/f_j)^2-\sum_I(w_i/u_i)^2+s^TLs\right\}=-\mu.
\]
This does not restrict intermediate powers in a Laurent parameter.
\end{proof}

\section{One order-26 support with opposite outcomes for \texorpdfstring{$\kappa$}{kappa}}
Set $\omega=e^{2\pi i/13}$, $\zeta=-\omega^7$, and $K=\Q(\omega)$. Take
\begin{align*}
 I&=\{0,1,2,3,5,7,13,15,17,18,19,20\},\\
 J&=\{3,9,13,14,15,16,18,20,22,23,24,25\},\\
 B=J^c&=\{0,1,2,4,5,6,7,8,10,11,12,17,19,21\}.
\end{align*}
For $E=\{0,1,2,3,4,5,6\}$ and $O=\{2,3,4,7,9,10,12\}$, write
\[
 U(X)=\prod_{e\in E}(X-\omega^e)=\sum a_iX^i,\qquad
 W(X)=\prod_{o\in O}(X+\omega^o)=\sum b_iX^i.
\]
Let $u$ be the length-$26$ coefficient vector of
\[
 b_1U(X)+a_1W(X)+X^{13}\{b_1U(X)-a_1W(X)\}.
\]
Exact arithmetic gives $\supp u=I$ and $\supp Fu=J$.
Put $D_0=I\setminus\{20\}$, $T_0=\{17,19,21\}$, $C_0=B\setminus T_0$.
The normal system uses
\begin{equation}\label{eq:Gexplicit}
 M=F[D_0,C_0],\qquad G_{T_0}=\mathrm{Id}_3,\qquad
 G_{C_0}=-M^{-1}F[D_0,T_0].
\end{equation}
Under $\zeta\mapsto5215$ modulo $10037$, $\det M=3759$.
Thus $\rank F[B,I]=11$, and $z=(g_{17},g_{19},g_{21})$ are free
coordinates. The modular minor proves nonvanishing only; exact cyclotomic
identities prove the support zeros.

\begin{theorem}\label{thm:26}
These supports and this $u$ admit algebraic nondegenerate projective
critical points of \eqref{eq:psi} with $\kappa\ne0$. They also admit
at least two distinct algebraic nondegenerate projective critical points with
$\kappa=0$. At two of the latter, $\mu\ne0$; hence no formal branch with
boundary direction $u$ and first transverse direction $w$, in the sense of
Theorem~\ref{thm:local}, leaves either of them.
\end{theorem}

\subsection{A point with \texorpdfstring{$\kappa\ne0$}{nonzero kappa}}
Reduce at the split prime $53$ with $\zeta\mapsto4$, $\omega\mapsto16$.
At $z=(1,8,19)$, exact finite-field arithmetic gives
\[
 \nabla\Psi(z)=0,\qquad
 H=\begin{pmatrix}6&47&5\\47&3&13\\5&13&11\end{pmatrix},\qquad\kappa=4.
\]
All $28$ denominators are nonzero, $\det M=46$, and the Hessian determinant
on $z_0=1$ is $23$. Hensel's lemma gives a normal critical point over
$\Q_{53}$ at which $\kappa$ is a unit. The cleared two-equation system has
an invertible Jacobian, so the point is isolated and algebraic over $K$.
Embed a number field of definition into $\C$ by an embedding extending
$\omega\mapsto e^{2\pi i/13}$. This gives the required point, preserving the
nonzero coordinates, the Hessian rank and $\kappa\ne0$. By \eqref{eq:kappaodd} below, $\kappa(u,Rz)=-\kappa(u,z)\ne0$. Since $\kappa$ scales as $\lambda^{-1}$ under $z\mapsto\lambda z$, $Rz$ could represent the same projective point only if $Rz=-z$, which forces $z_1=0$; but $z_1$ reduces to $8\ne0$ modulo $53$. So $Rz$ is a second such point.

\subsection{A symmetry that makes \texorpdfstring{$\kappa$}{kappa} odd}
Define monomial involutions, with indices modulo $26$, by
\begin{equation}\label{eq:involutions}
 (\mathcal Ux)_i=\zeta^{-10-12i}x_{20-i},\qquad
 (\mathcal Sg)_k=\zeta^{-10-20k}g_{12-k}.
\end{equation}
Direct multiplication gives
\begin{equation}\label{eq:intertwine}
 \mathcal U^2=\mathcal S^2=1,\qquad
 F\mathcal U=\mathcal S^{-T}F,\qquad
 F\mathcal S=\mathcal U^{-T}F.
\end{equation}
The supports are preserved since $I=20-I$ and $J=12-J$.
Substitution in the polynomial for $u$, in $\Z[\omega]$, gives
\begin{equation}\label{eq:uodd}
 \mathcal Uu=-u.
\end{equation}
Equivalently, $(\mathcal U_0x)_i=\zeta^{-12i}x_{20-i}$ acts on $u$
by $-\zeta^{10}$. On the free coordinates, $\mathcal S$ induces
\begin{equation}\label{eq:R}
 R(z_0,z_1,z_2)=(\zeta^{14}z_2,z_1,\zeta^{12}z_0).
\end{equation}
The products in \eqref{eq:psi} are permuted and multiplied by constants,
so $R$ preserves the critical locus. If $(u,g,w)$ satisfies \eqref{eq:w},
then $(\mathcal Uu,\mathcal Sg,\mathcal Uw)$ does also. Moreover
$h(\mathcal Uu)=\mathcal U^{-T}h(u)$, by \eqref{eq:intertwine}.
Therefore $\kappa(\mathcal Uu,\mathcal Sg)=\kappa(u,g)$.
Equation \eqref{eq:uodd} and scaling in $u$ imply
\begin{equation}\label{eq:kappaodd}
 \kappa(u,Rz)=-\kappa(u,z)
\end{equation}
on the normal critical locus. An exactly $R$-fixed representative has
$\kappa=0$. Projective fixedness without control of scale would not alone
justify that conclusion.

\subsection{Two fixed Hensel lifts}
Let $\pi=1-\omega$, with $\nu_\pi(\pi)=1$. The basis
\eqref{eq:Gexplicit} is integral at $\pi$, with the free rows the identity.
Divide each row of $G,V$ by its largest common power of $\pi$ and reduce.
Up to nonzero row constants, the $G$ rows, in increasing $B$ order, are
\begin{equation}\label{eq:rows}
\begin{gathered}
(1,4,7),(1,1,6),(1,3,9),(1,0,0),(1,7,4),(0,1,0),\\
(1,5,10),(0,0,1),(1,9,3),(1,11,11),(1,8,2),\\
(1,0,0),(0,1,0),(0,0,1).
\end{gathered}
\end{equation}
All $G$ rows have valuation zero. Every normalized $V$ row is proportional
to $\ell=(1,6,1)$; the row valuation is $4$ at the nine even indices
of $I^c$ and $8$ at its five odd indices. The accompanying certificate
program verifies these finite cyclotomic arithmetic statements. In particular,
write each entry as $\sum_{j=0}^{11}c_j\pi^j$, $c_j\in\Q$. Its valuation is
$\min_{c_j\ne0}(12v_{13}(c_j)+j)$: the term valuations are distinct modulo
$12$, so the minimum cannot cancel. Initial units follow from
$13/\pi^{12}\equiv-1\pmod\pi$. Thus these reductions can be checked in
the exact number field, without a precision cutoff. They give
the reduced logarithmic derivative of
\[
 \overline\Psi=\sum_{r\text{ in }\eqref{eq:rows}}\log(rz)
                   -14\log(z_0+6z_1+z_2).
\]
In $\F_{169}=\F_{13}[t]/(t^2-2)$ take
\[
 \bar z_\pm=(9\pm4t,\ 8\pm3t,\ 9\pm4t).
\]
They lie on $\ell z=1$, all $28$ normalized forms are nonzero, and the
gradient vanishes. On $z_2=1-z_0-6z_1$ the Hessian is
\[
 \begin{pmatrix}7\pm9t&8\pm t\\8\pm t&4\end{pmatrix},
 \qquad \det=1\pm7t\ne0,
\]
whose norm to $\F_{13}$ is $7$.

For the lift use the exactly $R$-invariant affine chart
\[
 L(z)=z_0+6z_1+\zeta^{14}z_2=1.
\]
This reduces to $\ell z=1$. On each residue disc the normalized
denominators are units. Hensel's lemma therefore gives a unique critical
lift on each disc in the unramified quadratic extension of $K_\pi$.
The reduction of $R$ swaps $z_0,z_2$ and fixes both residual points.
Since $R$ preserves $L=1$, the equations, and each disc, uniqueness forces
both lifts to be exactly $R$-fixed. Equation \eqref{eq:kappaodd} proves
$\kappa=0$ exactly. As above, nonsingularity of the polynomial system
makes both points algebraic over $K$. Embedding a common number field
of definition into $\C$ in the same way gives two distinct complex endpoints.

\subsection{The second scalar at the fixed lifts}
At each of the two lifts, $\kappa=0$, so $\mu$ is defined by \eqref{eq:mu}.
Normalize $u$ as above, $g=Gz$ in the chart $L(z)=1$, and $w_{20}=0$
(allowed since $u_{20}\ne0$); solve $Ls=-h$ with $s_4=s_{20}=0$. This is
possible because $\ker L=\operatorname{span}\{u,w\}$ and $u_4=0$ while $u_{20}w_4\ne0$,
so evaluation at the coordinates $4,20$ is injective on $\ker L$.
The accompanying program \texttt{certify\_mu.py} computes in
$\Z_{13}[\pi,t]$ with $t^2=2$ modulo $\pi^{240}$, carrying for every
quantity a valuation and a guaranteed absolute precision (conservative
$\pi$-adic ball arithmetic, no floating point). Newton's method in the chart,
started at $\bar z_\pm$, reaches gradient valuation $232$; since the
Hessian is a unit on the whole residue disc, the exact Hensel lift lies within
$O(\pi^{232})$ of the iterate, and all later quantities are computed on that
enclosure. It finds
\[
 v_\pi(\mu_\pm)=-14,\qquad \pi^{14}\mu_\pm\equiv 11\pm8t\pmod{\pi},
\]
with absolute precision $193$; the residues have norm $6\ne0$ in
$\F_{13}$. Thus $\mu\ne0$ at both lifts. Being nonzero is independent of
the normalizations, by the scaling law after \eqref{eq:mu}, and is preserved
by the embedding into $\C$, since $\mu$ is an algebraic number. By
Theorem~\ref{thm:local}, $Q_{\rm eff}'(\tau)=-\mu\tau+O(\tau^2)$ is not
identically zero, so no branch with the directions $u,w$ exists. This proves
Theorem~\ref{thm:26}. It excludes neither other directions through these
boundary points nor other strata.

\section{What remains open}
The multiset theorem supplies a necessary condition at every gap, without
classifying which support pairs can come from a single Laurent solution.
The local reduction is exact under its rank and nondegeneracy hypotheses.
Degenerate normal critical points, larger endpoint kernels, and branches
with more than two valuation levels require separate treatment.

Even in the nondegenerate case, $Q_{\rm eff}'\equiv0$ gives no bound
on how many coefficients must be examined. The order-$26$ symmetry proves
that testing $\kappa$ alone cannot exclude every endpoint; at the two
certified points $\mu$ does exclude the prescribed branches. Whether some
fixed number of coefficients always suffices is open. A uniform finiteness
argument needs an additional theorem forcing some coefficient to be
nonzero, with control over the endpoints and orders covered. The present
results do not supply that theorem.

\paragraph{Computational evidence and authorship.}
The certificates use exact cyclotomic arithmetic, finite fields and
conservative $\pi$-adic ball arithmetic.
A finite-precision zero is not used to prove $\kappa=0$; symmetry and
Hensel uniqueness prove that identity.
Two AI systems, ChatGPT/Codex (OpenAI) and Claude (Anthropic), worked on
this program in parallel under the author's direction and contributed to
finding the results, attacking and checking the proofs, verifying the
certificates, searching prior work and drafting the text. AI is not an author
or a referee, AI review is not peer review, and agreement between AI systems
is not independent validation. The author is responsible for the claims and
their presentation.

\begin{thebibliography}{9}
\bibitem{AV} D. Adrovic and J. Verschelde,
\emph{Polyhedral methods for space curves exploiting symmetry applied to the
cyclic $n$-roots problem}, in Computer Algebra in Scientific Computing (CASC 2013), Lecture Notes in Comput. Sci., Springer, 2013, 10--29, \url{https://doi.org/10.1007/978-3-319-02297-0_2}; arXiv:1109.0241.
\bibitem{Backelin} J. Backelin,
\emph{Square multiples $n$ give infinitely many cyclic $n$-roots},
Reports, Matematiska Institutionen, Stockholms Universitet, no.~8 (1989).
\bibitem{BS95} G. Bj\"orck and B. Saffari,
\emph{New classes of finite unimodular sequences with unimodular Fourier
transforms. Circulant Hadamard matrices with complex entries},
C. R. Acad. Sci. Paris S\'er.~I Math. \textbf{320} (1995), 319--324.
\bibitem{CDFV} D. Cohen, G. Denham, M. Falk, and A. Varchenko,
\emph{Critical points and resonance of hyperplane arrangements},
Canadian Journal of Mathematics \textbf{63} (2011), 1038--1057,
\url{https://doi.org/10.4153/CJM-2011-028-8}.
\bibitem{DP} C. De Concini and C. Procesi,
\emph{Wonderful models of subspace arrangements},
Selecta Math. (N.S.) \textbf{1} (1995), 459--494,
\url{https://doi.org/10.1007/BF01589496}.
\bibitem{Fei} E. M. Feichtner,
\emph{De Concini--Procesi wonderful arrangement models: a discrete
geometer's point of view},
\url{https://arxiv.org/abs/math/0403183}.
\bibitem{FS} E. M. Feichtner and B. Sturmfels,
\emph{Matroid polytopes, nested sets and Bergman fans},
\url{https://arxiv.org/abs/math/0411260}.
\bibitem{FR} H. F\"uhr and Z. Rzeszotnik,
\emph{On biunimodular vectors for unitary matrices},
Linear Algebra Appl. \textbf{484} (2015), 86--129, \url{https://doi.org/10.1016/j.laa.2015.06.019}; arXiv:1506.06738.
\bibitem{Haagerup} U. Haagerup,
\emph{Cyclic $p$-roots of prime length $p$ and related complex Hadamard matrices},
arXiv:0803.2629 (2008), \url{https://arxiv.org/abs/0803.2629}.
\bibitem{SvdV} A.-L. Sattelberger and R. van der Veer,
\emph{Maximum likelihood estimation from a tropical and a Bernstein--Sato
perspective}, Int. Math. Res. Not. IMRN \textbf{2023} (2023), no.~6,
5263--5292, \url{https://doi.org/10.1093/imrn/rnac016}; arXiv:2101.03570.
\bibitem{Zhou} W. Zhou,
\emph{On number of cyclic $n$-roots and disjointness of Fourier supports},
arXiv:2608.06335v2 (2026), \url{https://arxiv.org/abs/2608.06335v2}.
\end{thebibliography}
\end{document}
